\section{Introduction}
\label{sec:introduction}

% --- Motivation ---
The transmission of monetary policy shocks remains a central question of
quantitative macroeconomics. Because nominal prices adjust only sluggishly,
a contractionary innovation to the policy interest rate raises the real
rate, depresses aggregate demand, and feeds back into inflation through
the Phillips curve. The size and timing of these responses are
quantitatively important: they discipline the conduct of monetary policy,
they shape the welfare cost of inflation stabilisation, and they provide
the benchmark against which richer empirical dynamic stochastic general
equilibrium (DSGE) models are evaluated. The three-equation New Keynesian
model, despite its parsimony, captures the essential propagation
mechanism and remains the workhorse for teaching, policy communication,
and analytical reference.

% --- Literature ---
The closest references for the present paper are
\citet{KydlandPrescott1982}, who pioneered the modern quantitative
business-cycle research programme; \citet{Calvo1983}, whose
staggered-pricing assumption underpins the supply side; \citet{Taylor1993},
whose policy rule defines the interest-rate response to inflation and
output; \citet{Woodford2003} and \citet{Gali2008}, who provide the
canonical book-length treatments of the New Keynesian framework, including
the parameter values and impulse responses that I follow; \citet{Klein2000}
and \citet{BlanchardKahn1980}, whose solution methods underpin the
numerical approach used here; and \citet{SmetsWouters2007}, whose
medium-scale extension demonstrates the empirical relevance of the
broader framework in a Bayesian estimation context. The present paper
does not extend these contributions theoretically; instead, it serves as
a transparent quantitative replication of the canonical New Keynesian
monetary transmission with explicit documentation of the steady state,
the policy function, and the impulse responses.

% --- Contribution ---
Three contributions follow. First, I provide a transparent solution of
the three-equation New Keynesian model using the generalised Schur
decomposition of \citet{Klein2000}, deriving the closed-form policy
matrix and verifying the Blanchard--Kahn condition numerically. Second,
I document three quantitative features of the impulse response to a 25
basis-point contractionary monetary policy shock: on-impact output falls
by approximately one-quarter of one percent, on-impact inflation falls
by approximately one-third of one percentage point on an annualised
basis, and the real interest rate rises by approximately one-half of one
percentage point on an annualised basis. Third, each reported number
admits a closed-form check against the steady state of the linearised
system, so that the link between calibrated parameters, policy
coefficients, and impulse responses is fully transparent.

% --- Roadmap ---
The remainder of the paper is organised as follows.
Section~\ref{sec:model} presents the model.
Section~\ref{sec:calibration} discusses the calibration and steady-state
properties.
Section~\ref{sec:solution} describes the solution method and verifies
stability.
Section~\ref{sec:results} reports and interprets the impulse responses.
Section~\ref{sec:robustness} examines robustness to the Taylor-rule
reaction coefficients.
Section~\ref{sec:conclusion} concludes.
